\begin{lemma}
\label{lemma-compute-ext}
Let $f : X \to S$ be a morphism of schemes, $E \in D(\mathcal{O}_X)$
and $\mathcal{G}^\bullet$ a complex of $\mathcal{O}_X$-modules.
Assume
\begin{enumerate}
\item $S$ is Noetherian,
\item $f$ is locally of finite type,
\item $E \in D^-_{\textit{Coh}}(\mathcal{O}_X)$,
\item $\mathcal{G}^\bullet$ is a bounded complex of
coherent $\mathcal{O}_X$-modules flat over $S$ with support proper over $S$.
\end{enumerate}
Then the following two statements are true
\begin{enumerate}
\item[(A)] for every $m \in \mathbf{Z}$ there exists a perfect object $K$
of $D(\mathcal{O}_S)$ and functorial maps
$$
\alpha^i_\mathcal{F} :
\Ext^i_{\mathcal{O}_X}(E,
\mathcal{G}^\bullet \otimes_{\mathcal{O}_X} f^*\mathcal{F})
\longrightarrow
H^i(S, K \otimes^\mathbf{L}_{\mathcal{O}_S} \mathcal{F})
$$
for $\mathcal{F}$ quasi-coherent on $S$ compatible with boundary maps
(see proof) such that $\alpha^i_\mathcal{F}$ is an isomorphism for $i \leq m$
\item[(B)] there exists a pseudo-coherent $L \in D(\mathcal{O}_S)$
and functorial isomorphisms
$$
\Ext^i_{\mathcal{O}_S}(L, \mathcal{F}) \longrightarrow
\Ext^i_{\mathcal{O}_X}(E,
\mathcal{G}^\bullet \otimes_{\mathcal{O}_X} f^*\mathcal{F})
$$
for $\mathcal{F}$ quasi-coherent on $S$ compatible with boundary maps.
\end{enumerate}
\end{lemma}

\begin{proof}
Proof of (A).
Suppose $\mathcal{G}^i$ is nonzero only for $i \in [a, b]$.
We may replace $X$ by a quasi-compact open neighbourhood of
the union of the supports of $\mathcal{G}^i$.
Hence we may assume $X$ is Noetherian.
In this case $X$ and $f$ are quasi-compact and quasi-separated.
Choose an approximation $P \to E$ by a perfect complex $P$ of
$(X, E, -m - 1 + a)$
(possible by Theorem \ref{theorem-approximation}).
Then the induced map
$$
\Ext^i_{\mathcal{O}_X}(E,
\mathcal{G}^\bullet \otimes_{\mathcal{O}_X} f^*\mathcal{F})
\longrightarrow
\Ext^i_{\mathcal{O}_X}(P,
\mathcal{G}^\bullet \otimes_{\mathcal{O}_X} f^*\mathcal{F})
$$
is an isomorphism for $i \leq m$. Namely, the kernel, resp.\ cokernel of this
map is a quotient, resp.\ submodule of
$$
\Ext^i_{\mathcal{O}_X}(C,
\mathcal{G}^\bullet \otimes_{\mathcal{O}_X} f^*\mathcal{F})
\quad\text{resp.}\quad
\Ext^{i + 1}_{\mathcal{O}_X}(C,
\mathcal{G}^\bullet \otimes_{\mathcal{O}_X} f^*\mathcal{F})
$$
where $C$ is the cone of $P \to E$. Since $C$ has vanishing cohomology
sheaves in degrees $\geq -m - 1 + a$ these $\Ext$-groups are zero
for $i \leq m + 1$ by
Derived Categories, Lemma \ref{derived-lemma-negative-exts}.
This reduces us to the case that
$E$ is a perfect complex which is Lemma \ref{lemma-compute-ext-perfect}.
The statement on boundaries is explained in the proof of
Lemma \ref{lemma-compute-ext-perfect}.

\medskip\noindent
Proof of (B). As in the proof of (A) we may assume $X$ is Noetherian.
Observe that $E$ is pseudo-coherent by
Lemma \ref{lemma-identify-pseudo-coherent-noetherian}.
By Lemma \ref{lemma-pseudo-coherent-hocolim} we can write
$E = \text{hocolim} E_n$ with $E_n$ perfect and $E_n \to E$ inducing
an isomorphism on truncations $\tau_{\geq -n}$. Let $E_n^\vee$
be the dual perfect complex
(Cohomology, Lemma \ref{cohomology-lemma-dual-perfect-complex}).
We obtain an inverse system $\ldots \to E_3^\vee \to E_2^\vee \to E_1^\vee$
of perfect objects. This in turn gives rise to an inverse system
$$
\ldots \to K_3 \to K_2 \to K_1\quad\text{with}\quad
K_n = Rf_*(E_n^\vee \otimes_{\mathcal{O}_X}^\mathbf{L} \mathcal{G}^\bullet)
$$
perfect on $S$, see Lemma \ref{lemma-tensor-perfect}.
By Lemma \ref{lemma-compute-ext-perfect} and its proof and
by the arguments in the previous paragraph (with $P = E_n$)
for any quasi-coherent $\mathcal{F}$ on $S$ we have
functorial canonical maps
$$
\xymatrix{
& \Ext^i_{\mathcal{O}_X}(E,
\mathcal{G}^\bullet \otimes_{\mathcal{O}_X} f^*\mathcal{F})
\ar[ld] \ar[rd] \\
H^i(S, K_{n + 1} \otimes_{\mathcal{O}_S}^\mathbf{L} \mathcal{F})
\ar[rr] & &
H^i(S, K_n \otimes_{\mathcal{O}_S}^\mathbf{L} \mathcal{F})
}
$$
which are isomorphisms for $i \leq n + a$.
Let $L_n = K_n^\vee$ be the dual perfect complex.
Then we see that $L_1 \to L_2 \to L_3 \to \ldots$
is a system of perfect objects in $D(\mathcal{O}_S)$
such that for any quasi-coherent $\mathcal{F}$ on $S$
the maps
$$
\Ext^i_{\mathcal{O}_S}(L_{n + 1}, \mathcal{F})
\longrightarrow
\Ext^i_{\mathcal{O}_S}(L_n, \mathcal{F})
$$
are isomorphisms for $i \leq n + a - 1$. This implies that
$L_n \to L_{n + 1}$ induces an isomorphism on truncations
$\tau_{\geq -n - a + 2}$ (hint: take cone of $L_n \to L_{n + 1}$
and look at its last nonvanishing cohomology sheaf).
Thus $L = \text{hocolim} L_n$ is pseudo-coherent, see
Lemma \ref{lemma-pseudo-coherent-hocolim}. The mapping property
of homotopy colimits gives that
$\Ext^i_{\mathcal{O}_S}(L, \mathcal{F}) =
\Ext^i_{\mathcal{O}_S}(L_n, \mathcal{F})$
for $i \leq n + a - 3$ which finishes the proof.
\end{proof}